\editionPage{progression}{Suivre ses progrès}{Un bilan utile après chaque chapitre}{Cocher après avoir travaillé sans le corrigé.}
\begin{tabularx}{\linewidth}{@{}rYccc@{}}\toprule
N° & Chapitre & Cours & Exos & Défi\\\midrule
01 & Mesurer et raisonner & \caseSuivi & \caseSuivi & \caseSuivi \\
02 & Identifier la matière & \caseSuivi & \caseSuivi & \caseSuivi \\
03 & Préparer et doser une solution & \caseSuivi & \caseSuivi & \caseSuivi \\
04 & Atomes, noyaux et ions & \caseSuivi & \caseSuivi & \caseSuivi \\
05 & Électrons et stabilité & \caseSuivi & \caseSuivi & \caseSuivi \\
06 & Compter les entités & \caseSuivi & \caseSuivi & \caseSuivi \\
07 & Changements d’état et énergie & \caseSuivi & \caseSuivi & \caseSuivi \\
08 & Réactions chimiques et bilans & \caseSuivi & \caseSuivi & \caseSuivi \\
09 & Synthétiser et identifier & \caseSuivi & \caseSuivi & \caseSuivi \\
10 & Transformations nucléaires & \caseSuivi & \caseSuivi & \caseSuivi \\
11 & Décrire un mouvement & \caseSuivi & \caseSuivi & \caseSuivi \\
12 & Forces et interactions & \caseSuivi & \caseSuivi & \caseSuivi \\
13 & Inertie et chute libre & \caseSuivi & \caseSuivi & \caseSuivi \\
14 & Émettre et transmettre un son & \caseSuivi & \caseSuivi & \caseSuivi \\
15 & Période et perception sonore & \caseSuivi & \caseSuivi & \caseSuivi \\
16 & Lumière et spectres & \caseSuivi & \caseSuivi & \caseSuivi \\
17 & Réflexion et réfraction & \caseSuivi & \caseSuivi & \caseSuivi \\
18 & Lentilles et vision & \caseSuivi & \caseSuivi & \caseSuivi \\
19 & Circuits et lois électriques & \caseSuivi & \caseSuivi & \caseSuivi \\
20 & Capteurs et outils numériques & \caseSuivi & \caseSuivi & \caseSuivi \\
\bottomrule\end{tabularx}\par\medskip
\small Une case cochée signifie « je sais expliquer et refaire ». Si une méthode reste fragile, noter un exercice à reprendre puis le refaire quelques jours plus tard. Les erreurs de conversion, de signe et de légende méritent une attention particulière.
\clearpage
